% ---- pg-557 （书页 545） ----

组内平方和(到中心的距离平方和)为

\hspace*{6em}类 $A$:\hfill 0

\hspace*{6em}类 $(BCD)$:\hfill $4+5+5=14$

等价地, 我们可以用准则
\[
\min E=\sum d^{2}_{i,\text{类}(i)}
\]
确定 $K=2$ 个类, 其中最小化是对所有 $K=2$ 个类的聚类, $d^{2}_{i,\text{类}(i)}$ 是从项目 $i$ 到其所在类的中心的距离平方.

在此例中, 总共有 7 个可能的 $K=2$ 个类的聚类:

\hspace*{4em}$A,(BCD)$

\hspace*{4em}$B,(ACD)$

\hspace*{4em}$C,(ABD)$

\hspace*{4em}$D,(ABC)$

\hspace*{4em}$(AB),(CD)$

\hspace*{4em}$(AC),(BD)$

\hspace*{4em}$(AD),(BC)$

\noindent 对于聚类 $\{A,(BCD)\}$:

\hspace*{6em}$A$:\hfill $d^{2}_{A,\text{类}(A)}=0$

\hspace*{6em}$(BCD)$:\hfill $d^{2}_{B,\text{类}(B)}+d^{2}_{C,\text{类}(C)}+d^{2}_{D,\text{类}(D)}=4+5+5=14$

因此 $\sum d^{2}_{i,\text{类}(i)}=0+14=14$.

\noindent 对于其他聚类, 不难验证

\begin{mvtable}
\begin{tabular}{lcccc}
& $B,(ACD)$ & & $\sum d^{2}_{i,\text{类}(i)}$ & $=48.7$\\[4pt]
& $C,(ABD)$ & & $\sum d^{2}_{i,\text{类}(i)}$ & $=27.7$\\[4pt]
& $D,(ABC)$ & & $\sum d^{2}_{i,\text{类}(i)}$ & $=31.3$\\[4pt]
& $(AB),(CD)$ & & $\sum d^{2}_{i,\text{类}(i)}$ & $=28$\\[4pt]
& $(AC),(BD)$ & & $\sum d^{2}_{i,\text{类}(i)}$ & $=27$\\[4pt]
& $(AD),(BC)$ & & $\sum d^{2}_{i,\text{类}(i)}$ & $=51.3$
\end{tabular}
\end{mvtable}

由于聚类 $\{A,(BCD)\}$ 有最小 $\sum d^{2}_{i,\text{类}(i)}$, 这是最终得到的划分. \hfill$\blacksquare$

为检查聚类结果的稳定性, 应该以新的初始分割重新启动算法. 一旦聚类被确定, 就可通过对项目表中诸项目的重新排列来对聚类作直观解释:
使处于第一个聚类中的项目排在最前面, 处于第二个聚类中的项目排在下一个位置, 等等.
一份列有各聚类中心坐标和每个聚类内部方差的表同样有助于刻画组与组之间的差别.

\noindent\textbf{例 12.12}\quad (对公用事业公司进行聚类的 $K$ 均值法)

我们回到对公用事业公司进行聚类的问题, 其由数据由\hyperref[tab:12.4]{表 12.4} 中给出. 这次采用 $K$ 均值法, 并对 $K$ 取若干不同的值.
我们把 $K=4$ 和 $K=5$ 时分类结果汇总如下表. 一般来说, 特定 $K$ 值的选择并不是明确无误的, 这依赖于对分类对象的背景知识和基于数据的评估.
(基于数据的评估也许包括选择 $K$ 值, 以使组间变差相对组内变差达到极大. 有关的量度可以是 $|\mathbf{W}|/|\mathbf{B}+\mathbf{W}|$ [见\hyperref[eq:6-38]{式(6-38)}] 和 $\mathrm{tr}(\mathbf{W}^{-1}\mathbf{B})$.
